%%%%%%%%%%%%%%%%%%%%%%%%%%%%%%%%%%%%%%%%%%%%%%%%%%%%%%%%%%%%%%%%%%%%%%%
%% ROUND 3: PEDAGOGICAL ADDITIONS
%% Complete LaTeX for all outstanding interactive and explanatory elements.
%% Each piece specifies its placement location in the main lecture_notes.tex.
%%%%%%%%%%%%%%%%%%%%%%%%%%%%%%%%%%%%%%%%%%%%%%%%%%%%%%%%%%%%%%%%%%%%%%%


%======================================================================
% PART 1: NEW TCOLORBOX DEFINITION — "Stop and Think" boxes
%======================================================================
%
% PLACEMENT: After line 60 (after the \metaphor box definition),
%            in the "Custom boxes" preamble section.

\newtcolorbox{stopandthink}{
  colback=colspace!8, colframe=colspace!80!black,
  fonttitle=\bfseries, title={\raisebox{-0.1em}{\Large ?}\;\;Stop and Think},
  arc=2mm, boxrule=0.8pt
}

% Also define a misconception box (red-themed, distinct from the generic warning):
\newtcolorbox{misconception}{
  colback=residual!8, colframe=residual!80!black,
  fonttitle=\bfseries, title={Misconception Alert},
  arc=2mm, boxrule=0.8pt
}


%======================================================================
% PART 2: SEVEN "STOP AND THINK" BOXES
%======================================================================

% ─── STOP-AND-THINK 1 ─────────────────────────────────────────────────
% PLACEMENT: After line 151 (after the Projection Theorem, before
%            Section 2 "Where Coefficients Come From").

\begin{stopandthink}
\textbf{Exercise.}
Take $\by = (3, 1)^\top$ and $\bm{x} = (2, 1)^\top$ in $\R^2$.
The column space $\Col(\bm{x})$ is the line through the origin in the
direction $(2,1)$.
\begin{enumerate}[nosep]
  \item Sketch the vectors $\by$ and $\bm{x}$ in $\R^2$.
  \item Compute the projection of $\by$ onto $\Col(\bm{x})$ using the
        formula $\yhat = \bm{x}(\bm{x}^\top\bm{x})^{-1}\bm{x}^\top\by$.
        Where does the projection land on the line?
  \item Draw the residual vector $\be = \by - \yhat$.
        Verify visually that it is perpendicular to $\bm{x}$.
  \item Compute $\bm{x}^\top \be$ and confirm it equals zero.
\end{enumerate}
\emph{Answer check:} You should get $\yhat = \tfrac{7}{5}(2, 1)^\top
= (2.8,\, 1.4)^\top$ and $\be = (0.2,\, -0.4)^\top$.
\end{stopandthink}


% ─── STOP-AND-THINK 2 ─────────────────────────────────────────────────
% PLACEMENT: After line 264 (after the proof of hat matrix properties,
%            before "Geometric meaning" paragraph).

\begin{stopandthink}
\textbf{Quick calculation.}
Suppose you fit a model with $p = 2$ predictors (including the intercept)
to $n = 10$ observations.
\begin{enumerate}[nosep]
  \item What is $\tr(\bH)$?
  \item What is the average leverage $\bar{h} = \frac{1}{n}\sum_i h_{ii}$?
  \item If one observation has leverage $h_{ii} = 0.6$, how does this
        compare to the ``high leverage'' threshold $2p/n$?
        Should you investigate this point?
\end{enumerate}
\emph{Answers:} (1)~$\tr(\bH) = p = 2$. \quad
(2)~$\bar{h} = p/n = 0.2$. \quad
(3)~Threshold $= 2 \times 2 / 10 = 0.4$. Since $0.6 > 0.4$, yes---this
observation has three times the average leverage and deserves scrutiny.
\end{stopandthink}


% ─── STOP-AND-THINK 3 ─────────────────────────────────────────────────
% PLACEMENT: After line 338 (after "the right triangle
%            $(\tilde{\by},\;\tilde{\yhat},\;\be)$ in $\R^n$",
%            just before the TikZ figure).

\begin{stopandthink}
\textbf{Thought experiment: adding a random noise column.}

Suppose your current model has $R^2 = 0.60$.
You add a new column $\bm{z}$ to $\bX$, where $\bm{z}$ is pure random noise
(independent of everything).
\begin{enumerate}[nosep]
  \item Will SSE increase, decrease, or stay the same?
        \emph{Hint:} adding a column enlarges $\Col(\bX)$.
  \item Will SSR increase, decrease, or stay the same?
  \item Will $R^2$ increase, decrease, or stay the same?
  \item If you keep adding random noise columns until $p = n$,
        what happens to $R^2$?
\end{enumerate}
\emph{Answers:}
(1)~SSE can only decrease (or stay the same)---the projection onto a
larger subspace is always at least as close.
(2)~Since SST is fixed and SSE decreases, SSR must increase.
(3)~$R^2$ increases---even though the new column is pure noise!
(4)~When $p = n$, $\Col(\bX) = \R^n$, so $\yhat = \by$, $\be = \bm{0}$,
and $R^2 = 1$. The model ``explains'' everything but predicts nothing.
This is why $R^2$ alone is not a reliable measure of model quality.
\end{stopandthink}


% ─── STOP-AND-THINK 4 ─────────────────────────────────────────────────
% PLACEMENT: After line 517 (after the FWL Key Insight box,
%            before the "two-step projection" paragraph).

\begin{stopandthink}
\textbf{Two limiting cases of FWL.}

Recall that $\bhat_2 = (\bX_2^\top \bM_1 \bX_2)^{-1}\bX_2^\top \bM_1 \by$,
where $\bM_1 = \bI - \bH_1$ removes the influence of $\bX_1$.
\begin{enumerate}[nosep]
  \item \textbf{Case 1:} $\bX_1 \perp \bX_2$
    (i.e.\ $\bX_1^\top \bX_2 = \bm{0}$).
    What does $\bM_1 \bX_2$ simplify to?
    What does this tell you about the FWL coefficient $\bhat_2$
    compared to the coefficient from a simple regression of $\by$ on $\bX_2$ alone?
  \item \textbf{Case 2:} $\bX_2 \in \Col(\bX_1)$
    (i.e.\ $\bX_2$ is a linear combination of the columns of $\bX_1$).
    What does $\bM_1 \bX_2$ simplify to?
    What goes wrong?
\end{enumerate}
\emph{Answers:}
(1)~When $\bX_1 \perp \bX_2$, we have $\bH_1 \bX_2 = \bm{0}$, so
$\bM_1 \bX_2 = \bX_2$. The residualised predictor equals the original
predictor---controlling for $\bX_1$ changes nothing.
The FWL coefficient equals the simple regression coefficient.
This is why uncorrelated predictors give the same coefficients whether
you include them together or separately.

(2)~When $\bX_2 \in \Col(\bX_1)$, we have $\bH_1 \bX_2 = \bX_2$, so
$\bM_1 \bX_2 = \bm{0}$. The residualised predictor is the zero vector.
The formula $(\bX_2^\top \bM_1 \bX_2)^{-1}$ requires inverting a zero matrix,
which is impossible.
This is perfect multicollinearity: $\bX_2$ contains no information
beyond what $\bX_1$ already provides.
\end{stopandthink}


% ─── STOP-AND-THINK 5 ─────────────────────────────────────────────────
% PLACEMENT: After line 668 (after the "Influence = leverage x surprise"
%            Key Insight box, before Section 6 "When the Geometry Fails").

\begin{stopandthink}
\textbf{Ranking influence.}

Consider four observations in a simple linear regression.
Rank them from \emph{smallest} to \emph{largest} Cook's distance $D_i$.
\begin{enumerate}[nosep, label=(\alph*)]
  \item \textbf{Central + on-trend:}
        predictor value near $\bar{x}$, response close to the fitted line.
  \item \textbf{Central + outlier:}
        predictor value near $\bar{x}$, response far from the fitted line.
  \item \textbf{Extreme + on-trend:}
        predictor value far from $\bar{x}$, response close to the fitted line.
  \item \textbf{Extreme + outlier:}
        predictor value far from $\bar{x}$, response far from the fitted line.
\end{enumerate}
\emph{Answer:}
\begin{itemize}[nosep]
  \item (a) has low leverage \emph{and} low surprise $\Rightarrow$ smallest $D_i$.
  \item (c) has high leverage but low surprise $\Rightarrow$ small $D_i$
        (it reinforces the trend).
  \item (b) has low leverage but high surprise $\Rightarrow$ moderate $D_i$
        (it wants to move the line but lacks the mechanical advantage).
  \item (d) has high leverage \emph{and} high surprise $\Rightarrow$ largest $D_i$
        (the dangerous case).
\end{itemize}
Ranking: $(a) < (c) < (b) < (d)$.
Note that (c) and (b) can swap depending on the magnitudes,
but the extremes are always (a) smallest and (d) largest.
\end{stopandthink}


% ─── STOP-AND-THINK 6 ─────────────────────────────────────────────────
% PLACEMENT: After line 823 (after the Ridge "selective geometric surgery"
%            Key Insight box, before "Geometric interpretation" paragraph).

\begin{stopandthink}
\textbf{Ridge limits.}

Recall that $\bhat_{\text{ridge}} = (\bX^\top\bX + \lambda\bI)^{-1}\bX^\top\by$,
and in the eigendecomposition each component is shrunk by
$\lambda_k / (\lambda_k + \lambda)$.
\begin{enumerate}[nosep]
  \item What happens as $\lambda \to 0$?
        What does each shrinkage factor approach?
  \item What happens as $\lambda \to \infty$?
        What does each shrinkage factor approach?
        What does $\bhat_{\text{ridge}}$ approach?
  \item Is there a value of $\lambda$ for which Ridge is
        equivalent to ignoring the data entirely?
\end{enumerate}
\emph{Answers:}
(1)~As $\lambda \to 0$, each factor $\lambda_k/(\lambda_k + \lambda) \to 1$,
so $\bhat_{\text{ridge}} \to \bhat_{\text{OLS}}$.
Ridge with no penalty is OLS.

(2)~As $\lambda \to \infty$, each factor $\lambda_k/(\lambda_k + \lambda) \to 0$,
so $\bhat_{\text{ridge}} \to \bm{0}$.
Infinite penalty shrinks all coefficients to zero---the model predicts
$\hat{y}_i = 0$ for everyone (or $\hat{y}_i = \bar{y}$ if the intercept
is not penalised).

(3)~Yes: $\lambda = \infty$ means the model ignores the data completely
and defaults to the prior that all coefficients are zero.
Ridge traces a continuous path from ``trust the data fully'' ($\lambda=0$)
to ``trust the prior fully'' ($\lambda=\infty$).
\end{stopandthink}


% ─── STOP-AND-THINK 7 ─────────────────────────────────────────────────
% PLACEMENT: After line 859 (after "touching a polytope with an
%            ellipsoid", before the Comparison subsection).

\begin{stopandthink}
\textbf{Why corners attract LASSO solutions.}

Grab a piece of paper and try this:
\begin{enumerate}[nosep]
  \item In $\R^2$, draw the $L^1$ diamond
        $\{(\beta_1, \beta_2) : |\beta_1| + |\beta_2| \le 1\}$.
        It has corners at $(1,0)$, $(-1,0)$, $(0,1)$, $(0,-1)$.
  \item Draw a family of concentric ellipses centred at some point
        outside the diamond (this represents the RSS contours centred
        at $\bhat_{\text{OLS}}$).
  \item Starting from the smallest ellipse, gradually expand it until
        it first touches the diamond. Where does it touch?
  \item Now replace the diamond with the $L^2$ circle
        $\{(\beta_1, \beta_2) : \beta_1^2 + \beta_2^2 \le 1\}$
        and repeat. Where does the first contact occur now?
\end{enumerate}
\emph{Key observation:}
The diamond's corners are ``pointy''---an expanding ellipse is far more
likely to first touch a corner (where one coordinate is exactly zero)
than a smooth edge.
The circle has no corners, so the first contact is almost always at a
point where \emph{both} coordinates are nonzero.
This is why LASSO produces sparse solutions (exact zeros) and Ridge does not.
The geometry of the constraint set, not the algebra, is what drives sparsity.
\end{stopandthink}


%======================================================================
% PART 3: FIVE MISCONCEPTION WARNING BOXES
%======================================================================

% ─── MISCONCEPTION 1 ──────────────────────────────────────────────────
% PLACEMENT: After line 292 (after the existing Warning about
%            orthogonality being by construction, before Section 4
%            "Pythagoras Runs a Regression").

\begin{misconception}
\textbf{``The projection validates the model.''}

A common student error: ``My residuals are uncorrelated with my
predictors, so the model must be correct.'' This confuses a
\emph{property of the method} with \emph{evidence for the model}.

The orthogonality $\bX^\top \be = \bm{0}$ holds for \emph{any}
design matrix $\bX$, regardless of whether $\bX$ contains the
right variables, whether the relationship is linear, or whether
the errors are well-behaved. You could regress stock returns on
the phases of the moon and the residuals would still be
perpendicular to the moon-phase column.

Residual-predictor orthogonality tells you that OLS did its job
(found the closest point). It says nothing about whether the
column space was the right target.
\end{misconception}


% ─── MISCONCEPTION 2 ──────────────────────────────────────────────────
% PLACEMENT: After line 406 (after the R^2 Key Insight box,
%            before Subsection "The Gauss--Markov theorem").

\begin{misconception}
\textbf{``$R^2 = 0.8$ means the model is good.''}

$R^2$ measures a geometric angle---nothing more.
A high $R^2$ does \emph{not} mean:
\begin{itemize}[nosep]
  \item The model is correctly specified (you might have the wrong
        functional form or omitted variables).
  \item The model will predict well out of sample (overfitting
        inflates $R^2$ on training data).
  \item The coefficients are meaningful (multicollinearity can give
        $R^2 = 0.99$ with wildly unstable individual coefficients).
\end{itemize}
Conversely, a low $R^2$ does \emph{not} mean the model is useless.
In fields with inherently noisy data (e.g.\ daily stock returns),
$R^2 = 0.03$ can represent an economically significant and
statistically reliable effect.

$R^2$ answers exactly one question: ``What fraction of the total
variability in $\by$ is captured by the projection onto $\Col(\bX)$?''
Whether that fraction is ``good enough'' depends on the scientific
context, not on any universal threshold.
\end{misconception}


% ─── MISCONCEPTION 3 ──────────────────────────────────────────────────
% PLACEMENT: After line 534 (after the FWL two-step geometric
%            explanation, before Subsection "The effect of predictor
%            correlation").

\begin{misconception}
\textbf{``Controlling for $=$ holding constant physically.''}

When we say ``the effect of education on income, controlling for age,''
students often imagine an experiment where we literally hold someone's
age fixed and change only their education.
This is \emph{not} what regression does.

``Controlling for'' in regression means a \emph{statistical} operation:
we partial out the linear association of age with both income and
education, and examine what remains.
This is the FWL residualisation.
It does not physically intervene on anyone's age---it merely
removes the component of variation that is linearly predictable
from age.

The distinction matters because:
\begin{itemize}[nosep]
  \item If the relationship between age and the other variables is
        nonlinear, the linear residualisation may not fully remove
        age's influence.
  \item If age causes education (e.g.\ older cohorts had fewer
        educational opportunities), ``controlling for'' age may
        absorb part of the causal effect you are trying to measure.
  \item The phrase ``holding constant'' suggests a causal
        intervention, but regression only identifies associations
        within the residualised space. Causal conclusions require
        additional assumptions (e.g.\ no unmeasured confounders).
\end{itemize}
\end{misconception}


% ─── MISCONCEPTION 4 ──────────────────────────────────────────────────
% PLACEMENT: After line 745 (after the VIF diagnostic paragraph,
%            before Subsection "Omitted variable bias").

\begin{misconception}
\textbf{``Multicollinearity makes the model bad.''}

Multicollinearity inflates standard errors of individual coefficients,
but it does \emph{not}:
\begin{itemize}[nosep]
  \item Bias the OLS estimates (they remain unbiased).
  \item Ruin the overall prediction ($\yhat$ and $R^2$ can be fine).
  \item Invalidate the $F$-test for the joint significance of a group
        of collinear predictors.
\end{itemize}
Multicollinearity is only a problem \emph{if you need to interpret
individual coefficients}.
If your goal is prediction, collinear predictors may cause no harm
at all---the instability in individual $\hat{\beta}_j$ averages out in
$\yhat = \bX\bhat$.

The geometric picture makes this clear: the \emph{projection} $\yhat$
onto the column space is perfectly stable.
What is unstable is the \emph{decomposition} of $\yhat$ into
contributions along nearly parallel basis vectors.
The shadow is sharp; it is the coordinates of the shadow that jitter.
\end{misconception}


% ─── MISCONCEPTION 5 ──────────────────────────────────────────────────
% PLACEMENT: After line 781 (after the OVB discussion ending with
%            "which lies beyond the geometric framework",
%            before Section 7 "Constraining the Shadow").

\begin{misconception}
\textbf{``Clean residual plots rule out omitted variable bias.''}

Students sometimes reason: ``I checked my residual plots---no patterns,
no heteroscedasticity, no outliers---so my model must be correctly
specified.'' This is dangerously wrong.

Omitted variable bias is \emph{invisible} to residual diagnostics.
Both the ``short'' regression (omitting $\bX_2$) and the ``long''
regression (including $\bX_2$) have perfectly orthogonal residuals and
can pass every residual plot test with flying colours.
The problem is not \emph{how} you projected but \emph{which column
space} you projected onto.

A real-world example: regressing health outcomes on exercise frequency
without controlling for socioeconomic status will produce clean residuals
but biased estimates, because wealthier people both exercise more
and have better health for other reasons.
The residual plot looks fine because OLS always produces residuals
orthogonal to the included predictors---by construction, not by
virtue of model correctness.

The only safeguards against OVB are domain knowledge, causal reasoning,
and sensitivity analysis---not statistical diagnostics.
\end{misconception}


%======================================================================
% PART 4: FOUR NEW ANALOGIES
%======================================================================

% ─── ANALOGY 1: MOVIE-SCREEN ANALOGY FOR FWL ──────────────────────────
% PLACEMENT: After line 534, right after the existing two-step geometric
%            explanation, as a new paragraph before "The effect of predictor
%            correlation" subsection. (In practice, insert just before line 536.)

\medskip
\noindent\textbf{The movie-screen analogy.}
Imagine you are in a cinema, watching a scene projected onto a large screen
(the column space $\Col(\bX_1, \bX_2)$).
The projector beam is the data vector $\by$.
Now suppose you want to see \emph{only} the contribution of the second
actor ($\bX_2$), after removing everything the first actor ($\bX_1$)
contributes to the scene.

The FWL procedure is like placing a filter over the projector that blocks
all wavelengths corresponding to $\bX_1$.
What remains on the screen---$\bM_1 \by$ projected onto $\bM_1 \bX_2$---shows
you the ``pure $\bX_2$ movie.''
The key insight is that this filtered projection recovers \emph{exactly}
the same coefficient $\bhat_2$ as watching the full unfiltered movie
and decomposing algebraically.
Filtering first, then projecting, gives the same answer as projecting
the whole thing and reading off the $\bX_2$-coordinate.


% ─── ANALOGY 2: TELESCOPE ANALOGY FOR REGULARISATION ──────────────────
% PLACEMENT: After line 878 (after the Ridge/LASSO comparison table,
%            after "Both trade a small amount of bias..." paragraph,
%            before Section 8 "Testing with Geometry").

\medskip
\noindent\textbf{The telescope analogy.}
Think of the OLS estimate as looking through a telescope on a windy night.
The telescope points in the right direction on average (unbiasedness),
but the wind makes it shake violently (high variance from
multicollinearity or small sample size).
Each snapshot is centered on the truth, but individual frames are blurry
and useless.

Regularisation is like tightening the telescope's mount.
A tighter mount introduces a small systematic error---the telescope
now points slightly away from the true target (bias).
But the shaking is dramatically reduced (lower variance), so each
individual frame is much sharper and more useful.

Ridge tightens the mount uniformly in all directions.
LASSO tightens some axes completely (locking them to zero) while
leaving others free---it is like a mount that can only swivel along
certain axes, automatically selecting which directions matter.

The bias-variance trade-off is the trade-off between pointing accuracy
and shake reduction.
The optimal penalty $\lambda$ is the tightness that minimises the
total blur (mean squared error = bias$^2$ + variance).


% ─── ANALOGY 3: SHADOW-SHARPNESS ANALOGY FOR R² ──────────────────────
% PLACEMENT: After line 406 (after the R^2 Key Insight box, and after
%            Misconception 2 above; or immediately after the Key Insight
%            box if misconception is placed elsewhere).

\medskip
\noindent\textbf{The shadow-sharpness analogy.}
Think of $R^2$ as measuring how \emph{sharp} a shadow is.
When the flashlight is nearly aligned with the data vector
($\theta \approx 0$), the shadow $\yhat$ is almost as long as the
object $\by$---the shadow is a crisp, faithful reproduction of the
object ($R^2 \approx 1$).

When the flashlight is nearly perpendicular to the data
($\theta \approx 90^\circ$), the shadow shrinks to almost nothing---a
tiny smudge that preserves none of the object's structure ($R^2 \approx 0$).

A shadow with $R^2 = 0.5$ (i.e.\ $\theta = 45^\circ$) preserves half
the object's ``energy'' (squared length).
The shadow exists, but it is a washed-out version of the original.
Whether that is good enough depends on what you need the shadow for.


% ─── ANALOGY 4: HEARING-TEST ANALOGY FOR F-TEST ──────────────────────
% PLACEMENT: After line 921 (after the F-test Key Insight box,
%            before the Remark about $F = t^2$).

\medskip
\noindent\textbf{The hearing-test analogy.}
The $F$-test is like an audiologist's hearing test.
You are in a noisy room (the residual noise $\norm{\be}^2/(n-p)$).
The audiologist plays a tone (the signal: $\norm{\yhat - \yhat_R}^2/q$)
and asks: ``Can you hear it above the background noise?''

If the tone is barely louder than the noise ($F \approx 1$), you cannot
distinguish signal from static---the extra predictors are not doing
anything detectable.
If the tone is much louder ($F \gg 1$), you hear it clearly---the extra
predictors have captured real structure.

The $p$-value is the probability that background noise alone would
produce a tone at least this loud.
A small $p$-value means: ``It is very unlikely that random noise
could mimic this signal; the extra predictors are probably real.''

The number of extra predictors $q$ matters because each additional
predictor is like an additional frequency channel.
If you listen on many channels ($q$ large), you are more likely to hear
\emph{something} by chance, so the test adjusts for this by dividing
the total signal by $q$.


%======================================================================
% PART 5: THREE PREREQUISITE REVIEW PARAGRAPHS
%======================================================================

% ─── PREREQUISITE 1: LOEWNER ORDERING ─────────────────────────────────
% PLACEMENT: After line 430 (after the Gauss--Markov proof,
%            before Subsection "R^2 monotonicity").

\medskip
\begin{remark}[Prerequisite: the Loewner ordering]
The Gauss--Markov theorem uses the notation
$\bm{A} \succeq \bm{B}$ (the \emph{Loewner ordering}).
This means $\bm{A} - \bm{B}$ is positive semi-definite, i.e.\
$\bm{v}^\top(\bm{A} - \bm{B})\bm{v} \ge 0$ for all $\bm{v}$.

For variance matrices, $\text{Var}(\btilde) \succeq \text{Var}(\bhat)$
means: in \emph{every direction} $\bm{v}$, the variance of the linear
combination $\bm{v}^\top\btilde$ is at least as large as the variance of
$\bm{v}^\top\bhat$.
It is a much stronger statement than saying ``the total variance is
larger''---it says the OLS estimator wins direction by direction.

A useful consequence: if $\bm{A} \succeq \bm{B} \succeq \bm{0}$, then
every eigenvalue of $\bm{A}$ is at least as large as the corresponding
eigenvalue of $\bm{B}$ (when both are expressed in the same basis).
In the Gauss--Markov context, the ``extra'' term $\sigma^2\bD\bD^\top$
is positive semi-definite (it is a product $\bD\bD^\top$), so the
Loewner inequality $\text{Var}(\btilde) \succeq \text{Var}(\bhat)$ is
immediate.
\end{remark}


% ─── PREREQUISITE 2: EIGENDECOMPOSITION OF SYMMETRIC MATRICES ─────────
% PLACEMENT: After line 708 (after "What multicollinearity means"
%            and before "The eigenvalue perspective" paragraph).
%            Specifically, before "The eigenvalue decomposition of
%            $\bX^\top\bX$ makes this precise."

\medskip
\begin{remark}[Prerequisite: eigendecomposition of symmetric matrices]
Before proceeding, let us recall a fundamental fact from linear algebra.

Every real symmetric matrix $\bm{A} \in \R^{p \times p}$ can be
decomposed as
\[
  \bm{A} = \bQ \bm{\Lambda} \bQ^\top,
\]
where $\bQ = [\bm{q}_1 \mid \cdots \mid \bm{q}_p]$ is an orthogonal
matrix ($\bQ^\top \bQ = \bI$) whose columns are the \emph{eigenvectors},
and $\bm{\Lambda} = \diag(\lambda_1, \dots, \lambda_p)$ contains the
corresponding \emph{eigenvalues}.

Key properties:
\begin{itemize}[nosep]
  \item All eigenvalues of a real symmetric matrix are \emph{real}
        (not complex).
  \item The eigenvectors can be chosen to be \emph{orthonormal}.
  \item If $\bm{A}$ is also positive semi-definite (like $\bX^\top\bX$),
        all eigenvalues are non-negative: $\lambda_k \ge 0$.
  \item The eigenvalues measure the ``stretch'' of the matrix in each
        eigenvector direction: $\bm{A}\bm{q}_k = \lambda_k \bm{q}_k$.
\end{itemize}

The inverse (when it exists) is simply
$\bm{A}^{-1} = \bQ \bm{\Lambda}^{-1} \bQ^\top
= \bQ \, \diag(1/\lambda_1, \dots, 1/\lambda_p) \, \bQ^\top$.
This makes it clear that small eigenvalues produce large entries in the
inverse---the source of variance inflation under multicollinearity.
\end{remark}


% ─── PREREQUISITE 3: WHERE THE GAUSSIAN ENTERS ────────────────────────
% PLACEMENT: After line 879 (after the Ridge/LASSO comparison and
%            telescope analogy, before Section 8 "Testing with Geometry:
%            The $F$-Test"). This prepares the reader for the inference
%            machinery that follows.

\medskip
\begin{remark}[Prerequisite: where the Gaussian enters]
Everything up to this point has been \emph{purely geometric}---no
distributional assumptions were needed.
The Projection Theorem, the hat matrix properties, Pythagoras, FWL,
leverage, influence, multicollinearity, OVB, and even regularisation
all follow from linear algebra alone.

To perform \emph{inference}---hypothesis tests, $p$-values, confidence
intervals---we need a probability model.
The standard assumption is:
\[
  \by = \bX\bm{\beta} + \bm{\varepsilon},
  \qquad
  \bm{\varepsilon} \sim \mathcal{N}(\bm{0},\, \sigma^2 \bI_n).
\]
That is, the errors are independent, identically distributed, and
Gaussian (normally distributed).

Why does the Gaussian matter?
\begin{itemize}[nosep]
  \item \textbf{Exactness:} Under Gaussianity, $\bhat$ is exactly
        $\mathcal{N}(\bm{\beta},\, \sigma^2(\bX^\top\bX)^{-1})$---not
        just approximately. This gives us exact $t$- and $F$-distributions
        for test statistics, valid in any sample size.
  \item \textbf{Without Gaussianity:} The Central Limit Theorem still
        ensures that $\bhat$ is \emph{approximately} normal in large samples.
        So inference is approximately valid even when the true error
        distribution is non-Gaussian---but only asymptotically.
  \item \textbf{Geometry of the Gaussian:} A Gaussian vector projected
        onto a subspace remains Gaussian. This is why the hat matrix
        preserves normality: $\yhat = \bH\by$ and
        $\be = \bM\by$ are both Gaussian, and they are
        \emph{independent} (because $\bH\bM = \bm{0}$).
        Independence of $\yhat$ and $\be$ is what makes the $F$-test work.
\end{itemize}
The next section uses these distributional facts to build hypothesis tests.
\end{remark}


%======================================================================
% PART 6: TRANSITION SENTENCES BETWEEN SECTIONS
%======================================================================

% Each transition is a short paragraph to be inserted at the END of
% a section, just before the next \section command.

% ─── TRANSITION: Section 1 → Section 2 ────────────────────────────────
% PLACEMENT: After line 151 (after "S = Col(X)", before line 153
%            which is the Section 2 heading). Place BEFORE Stop-and-Think 1.

\medskip
\noindent\textit{The Projection Theorem guarantees that a closest point
exists and is characterised by orthogonality.
But how do we actually \emph{compute} it?
The next section translates the geometric existence result into an
explicit formula---the normal equations---and shows that the OLS
coefficients are simply the coordinates of this closest point.}


% ─── TRANSITION: Section 2 → Section 3 ────────────────────────────────
% PLACEMENT: After line 227 (end of Section 2 "Building intuition"
%            subsection, before line 230 which is the Section 3 heading).

\medskip
\noindent\textit{We now have a formula for the projection and its
coordinates.
But the formula hides a remarkable structure: the projection is a
\emph{linear operator} $\bH$ with powerful algebraic properties.
What can these properties tell us about individual observations,
and about the projection itself?}


% ─── TRANSITION: Section 3 → Section 4 ────────────────────────────────
% PLACEMENT: After line 293 (end of Section 3, after the Warning box
%            and Misconception 1 above, before Section 4 heading).

\medskip
\noindent\textit{The hat matrix decomposes $\by$ into two perpendicular
pieces: the shadow $\yhat$ and the residual $\be$.
Since these pieces are perpendicular, the Pythagorean theorem must
apply.
What does Pythagoras tell us about regression?
The answer is the variance decomposition---and $R^2$ as an angle.}


% ─── TRANSITION: Section 4 → Section 5 ────────────────────────────────
% PLACEMENT: After line 451 (end of Section 4, after the $R^2$
%            monotonicity Warning box, before line 454 which is the
%            Section 5 heading).

\medskip
\noindent\textit{So far, we have treated the column space as a monolithic
object and projected $\by$ onto it in one step.
But in multiple regression, different predictors contribute different
pieces of information.
Can we tease apart the contributions?
What does it really mean to ``control for'' a variable?
The Frisch--Waugh--Lovell theorem answers this by decomposing the
single projection into sequential steps.}


% ─── TRANSITION: Section 5 → Section 6 ────────────────────────────────
% PLACEMENT: After line 577 (end of Section 5, after the SE formula
%            and multicollinearity explanation, before line 580 which
%            is the Section 6 heading).

\medskip
\noindent\textit{FWL revealed that individual coefficients depend on
the unique information each predictor contributes, as measured by
the residualised predictor's length.
But some observations contribute more to the projection than others.
Which observations have the most influence on the fitted values,
and how can we detect problematic data points?}


% ─── TRANSITION: Section 6 → Section 7 ────────────────────────────────
% PLACEMENT: After line 781 (end of Section 6, after the OVB discussion
%            and Misconception 5 above, before Section 7 heading).

\medskip
\noindent\textit{Multicollinearity makes the column space geometrically
thin, inflating coefficient variance.
Omitted variable bias makes the column space the wrong target entirely.
The first problem has a geometric remedy: can we \emph{constrain}
the projection to trade a little bias for a lot of stability?
This is the idea behind regularisation.}


% ─── TRANSITION: Section 7 → Section 8 ────────────────────────────────
% PLACEMENT: After line 879 (end of Section 7, after the comparison
%            table and telescope analogy, after the Gaussian prerequisite
%            remark, before Section 8 heading).

\medskip
\noindent\textit{Regularisation reshapes the geometry to stabilise
estimates, but it cannot tell us whether the extra predictors are
\emph{statistically significant}.
For that, we need to compare two projections---one with and one
without the extra variables---and ask whether the improvement is
larger than what noise alone would produce.
This brings us to the $F$-test.}


% ─── TRANSITION: Section 8 → Section 9 ────────────────────────────────
% PLACEMENT: After line 945 (end of Section 8, after the confidence
%            ellipsoid paragraph "everything connects", before
%            Section 9 heading).

\medskip
\noindent\textit{The $F$-test and confidence ellipsoids complete our
inferential toolkit.
But before we summarise, it is worth stepping back and asking:
where does the geometric framework \emph{break down}?
What can projections not tell us, and where must we look beyond
linear algebra for answers?}


%======================================================================
% END OF ROUND 3 PEDAGOGICAL ADDITIONS
%======================================================================
